\documentclass[11pt]{article}
\usepackage[a4paper,margin=1in]{geometry}
\usepackage{amsmath,amssymb,amsthm}
\usepackage[T1]{fontenc}
\usepackage{lmodern}
\usepackage{microtype}
\usepackage{xcolor}
\usepackage{hyperref}

\hypersetup{
  colorlinks=true,
  linkcolor=blue!50!black,
  urlcolor=blue!50!black
}

\setlength{\parindent}{0pt}
\setlength{\parskip}{0.5em}

\newtheorem*{claim}{Claim}

\newcommand{\promptbox}[1]{%
  \begingroup\setlength{\fboxsep}{10pt}\noindent
  \colorbox{blue!8}{\parbox{\dimexpr\linewidth-2\fboxsep\relax}{\textbf{Prompt. }#1}}
  \endgroup
}
\newcommand{\resultbox}[1]{%
  \begingroup\setlength{\fboxsep}{10pt}\noindent
  \colorbox{purple!8}{\parbox{\dimexpr\linewidth-2\fboxsep\relax}{\textbf{Result. }#1}}
  \endgroup
}
\newcommand{\examplebox}[1]{%
  \begingroup\setlength{\fboxsep}{10pt}\noindent
  \colorbox{green!8}{\parbox{\dimexpr\linewidth-2\fboxsep\relax}{\textbf{Example. }#1}}
  \endgroup
}

\title{\textbf{Book 1 --- Lecture 1 --- Interlude 1 --- Exercise 3}\\[2mm]
\large Magnitude Squared and Dot Product}
\author{}
\date{}

\begin{document}
\maketitle

\promptbox{Show that the magnitude of a vector satisfies
\[
\lVert\mathbf{A}\rVert^2 = \mathbf{A}\cdot\mathbf{A}.
\]}

Let
\[
\mathbf{A} = (A_1, A_2, \dots, A_n).
\]

\begin{claim}
\[
\lVert\mathbf{A}\rVert^2 = \mathbf{A}\cdot\mathbf{A}.
\]
\end{claim}

\begin{proof}
By definition of Euclidean magnitude,
\[
\lVert\mathbf{A}\rVert = \sqrt{A_1^2 + A_2^2 + \cdots + A_n^2}.
\]
Squaring both sides gives
\[
\lVert\mathbf{A}\rVert^2 = A_1^2 + A_2^2 + \cdots + A_n^2.
\]

By definition of dot product,
\[
\mathbf{A}\cdot\mathbf{A}
= (A_1,\dots,A_n)\cdot(A_1,\dots,A_n)
= A_1^2 + A_2^2 + \cdots + A_n^2.
\]
Hence,
\[
\lVert\mathbf{A}\rVert^2 = \mathbf{A}\cdot\mathbf{A}. \qedhere
\]
\end{proof}

\resultbox{The identity follows directly from component definitions of the
Euclidean norm and dot product: both reduce to the same sum of squares.}

\examplebox{For \(\mathbf{A}=(3,-4,12)\),
\[
\mathbf{A}\cdot\mathbf{A} = 3^2 + (-4)^2 + 12^2 = 169,\qquad
\lVert\mathbf{A}\rVert = 13,\qquad
\lVert\mathbf{A}\rVert^2 = 169.
\]
So the identity holds numerically.}

\end{document}
